\documentclass[12pt]{article}
\usepackage[utf8]{inputenc}
\usepackage[T1]{fontenc}
\usepackage{hyperref}
\title{Ewondo Language Resources for K-12 Curriculum Development}
\author{Compiled by Claude Sonnet 5.5 (claude-sonnet-5-5) for Andr\'as Kornai}
\date{}
\begin{document}
\maketitle

\section{Introduction}
Ewondo (ISO 639-3 \texttt{ewo}) is a Bantu language of central Cameroon with roughly 580,000 speakers. This report lists only what a web search turned up. Little pedagogical material was found; most items are descriptive or encyclopedic.

\section{Primary Educational Resources}
No Ewondo primers, textbooks or school curricula were found in the search results. This is a gap for curriculum development.

\section{Linguistic and Reference Materials}
\begin{itemize}
\item James E. Redden, \emph{A Descriptive Grammar of Ewondo} (1979), listed in Glottolog: \url{https://glottolog.org/resource/reference/id/146626}
\item Afranaph bibliography of Ewondo (lists grammars, theses and studies, including work on harmonised orthography): \url{https://afranaphproject.afranaphdatabase.com/images/bibliographies/Ewondo__Afranaph-Bibliography.pdf}
\item Wikipedia overview of the language: \url{https://en.wikipedia.org/wiki/Ewondo_language}
\item Omniglot page (writing system, tones): \url{https://mail.omniglot.com/writing/ewondo.php}
\item Lexvo/CoCoon (Huma-Num) language entry: \url{https://cocoon.huma-num.fr/exist/crdo/lexvo.xql?code=ewo}
\end{itemize}

\section{Religious and Community Resources}
No Ewondo Bible translation or community organisation was identified in the search results.

\section{Digital and Audio Resources}
\begin{itemize}
\item Glosbe community-built Ewondo word translations (user-contributed, quality unverified): \url{https://glosbe.com/ewo/lib}
\end{itemize}
No audio resources were found.

\section{Curriculum Development Context}
The available material is mostly 1970s--80s descriptive linguistics, which can inform orthography and grammar sections but is not classroom-ready. Further searching of library catalogues is advisable.

\end{document}